\documentclass{article}

\usepackage[a4paper, margin=1in]{geometry}
\usepackage{listings}
\usepackage{xcolor}

\begin{document}

\newpage

\section*{Lab Report 1: GCD (Euclidean Algorithm)}

\subsection*{Theory}
The Greatest Common Divisor (GCD) of two integers is the largest positive integer that divides both numbers without leaving a remainder. The most efficient classical algorithm to compute GCD is the Euclidean Algorithm, based on the principle that GCD(a, b) = GCD(b, a mod b), repeating until the remainder becomes zero.

The iterative approach avoids recursion overhead, using a simple loop to repeatedly replace (a, b) with (b, a \% b) until b becomes 0, at which point a holds the GCD. This algorithm has deep roots in number theory and is one of the oldest known algorithms, attributed to Euclid (~300 BC).

\subsection*{Algorithm (Pseudocode)}
\begin{verbatim}
    ALGORITHM GCD_Iterative(a, b)
        Input: Two non-negative integers a and b
        Output: GCD of a and b

        WHILE b != 0 DO
            temp <- b
            b <- a MOD b
            a <- temp
        END WHILE
        RETURN a
\end{verbatim}

\subsection*{Python Program}
\begin{verbatim}
    def gcd_iterative(a, b):
        while b != 0:
            a, b = b, a % b

        return a


    if __name__ == "__main__":
        test_cases = [(48, 18), (100, 75), (56, 98), (0, 5), (17, 13)]

        print(f"{'a':>5} {'b':>5} {'GCD':>5}")
        print("-" * 21)

        for a, b in test_cases:
            print(f"{a:>5} {b:>5} {gcd_iterative(a, b):>5}")
\end{verbatim}

\subsection*{Output}
\begin{verbatim}
    a     b    GCD
---------------------
   48    18     6
  100    75    25
   56    98    14
    0     5     5
   17    13     1
\end{verbatim}

\subsection*{Complexity Analysis}

\subsubsection*{Time Complexity}

\begin{enumerate}
    \item Best Case: $O(1)$

          The best case occurs when the algorithm terminates in a single step i.e., either $b = 0$ or $a \: mod \: b = 0$

          Thus, $B(n) = O(1)$

    \item Average Case: $O(log\:n)$

          Each step replaces $(a,b) \rightarrow  (b,a\:mod\:b)$

          The remainder is significantly smaller than the previous number which leads to a logarithmic number of steps relative to the size of the input.

          Thus, $A(n) = O(log\:n)$

    \item Worst Case: $O(log\:min(a, b))$

          The worst case occurs when the reduction is as slow as possible. This happens when a and b are consecutive Fibonacci numbers.

          Thus, $W(n) = O(log\:min(a, b))$
\end{enumerate}

\subsubsection*{Space Complexity}
Only a few variables used $(a, b, r)$ for the iterative method.

Thus, $S(n) = O(1)$

\subsection*{Conclusion}
The iterative GCD algorithm using the Euclidean method is highly efficient. Its logarithmic time complexity $O(log\:min(a,b))$ makes it scalable even for very large integers. The constant $O(1)$ space usage is optimal. It is the standard algorithm used in cryptographic applications (RSA key generation, modular arithmetic) due to its speed and correctness.
\end{document}